\subsection{平坦模}

命题 1. 设 E 是右 A-模。下列三个性质等价：

\begin{enumerate}
\def\labelenumi{(\alph{enumi})}
\item
  E 对 \(A_{s}\) 是平坦的（换句话说，对 A 的每个左理想 a，典范同态 \(E \otimes_{A} a \to E \otimes_{A} A_{s} = E\) 是单射的）。
\item
  E 对每个左 A-模 M 都是 M-平坦的。
\item
  对左 A-模与同态的每个正合列
\end{enumerate}

\[
\mathrm{M} ^ {\prime} \xrightarrow {v} \mathrm{M} \xrightarrow {w} \mathrm{M} ^ {\prime \prime}
\]

序列

\[
\mathrm{E} \otimes_ {\mathrm{A}} \mathrm{M} ^ {\prime} \xrightarrow {1 \otimes v} \mathrm{E} \otimes_ {\mathrm{A}} \mathrm{M} \xrightarrow {1 \otimes w} \mathrm{E} \otimes_ {\mathrm{A}} \mathrm{M} ^ {\prime \prime}
\]

是正合的。

(b) 蕴含 (a) 是直接的。反之，设 (a) 成立；由第 2 目引理 5，E 对每个自由左 A-模都是平坦的；由于每个左 A-模都同构于一个自由模的商（《代数学》第二章 §1 第 11 目命题 20），由第 2 目的引理 4 即得 E 对 M 是平坦的。

我们证明 (c) 蕴含 (b)。若 \(v \colon \mathbf{M}' \to \mathbf{M}\) 是单同态，则序列 \(0 \to \mathbf{M}' \xrightarrow{v} \mathbf{M}\) 正合；由 (c)，序列

\[
0 \rightarrow \mathrm{E} \otimes_ {\mathrm{A}} \mathrm{M} ^ {\prime} \xrightarrow {1 \otimes v} \mathrm{E} \otimes_ {\mathrm{A}} \mathrm{M}
\]

是正合的；这表明 \(1 \otimes v\) 是单射，换句话说，E 是 M-平坦的。

最后，蕴含关系 \((\mathbf{b}) \Rightarrow (\mathbf{c})\) 是下面更精确的引理的推论：

引理 6. 若 \(\mathbf{M}' \xrightarrow{v} \mathbf{M} \xrightarrow{w} \mathbf{M}''\) 是左 A-模的一个正合列，且 E 是 \(\mathbf{M}''\) -平坦的右 A-模，则序列

\[
\mathrm{E} \otimes_ {\mathrm{A}} \mathrm{M} ^ {\prime} \xrightarrow {1 \otimes v} \mathrm{E} \otimes_ {\mathrm{A}} \mathrm{M} \xrightarrow {1 \otimes w} \mathrm{E} \otimes_ {\mathrm{A}} \mathrm{M} ^ {\prime \prime}
\]

是正合的。

我们以证明第 2 目引理 4 时的同一含义使用记号 \(\mathrm{T}(\mathbf{Q})\) 和 \(\mathrm{T}(v)\) 。记 \(\mathbf{M}_1'' = w(\mathbf{M})\) ，用 \(i: \mathbf{M}_1'' \to \mathbf{M}''\) 表示典范单射，用 \(p\) 表示 \(\mathbf{M}\) 到 \(\mathbf{M}_1''\) 的与 \(w\) 有相同图的映射。序列 \(\mathbf{M}' \xrightarrow{v} \mathbf{M} \xrightarrow{p} \mathbf{M}_1'' \to 0\) 正合，且由第 1 目引理 1 得序列

\[
\mathrm{T} (\mathbf {M} ^ {\prime}) \xrightarrow {\mathrm{T} (v)} \mathrm{T} (\mathbf {M}) \xrightarrow {\mathrm{T} (p)} \mathrm{T} (\mathbf {M} _ {1} ^ {\prime \prime}) \longrightarrow 0
\]

是正合的。此外，由于 E 对 \(\mathbf{M}''\) 是平坦的，映射 \(\mathrm{T}(i): \mathrm{T}(\mathbf{M}_1'') \to \mathrm{T}(\mathbf{M}'')\) 是单射的，且由于 \(\mathrm{T}(i) \circ \mathrm{T}(p) = \mathrm{T}(w)\) ，序列

\[
\mathrm{T} (\mathbf {M} ^ {\prime}) \xrightarrow {\mathrm{T} (v)} \mathrm{T} (\mathbf {M}) \xrightarrow {\mathrm{T} (w)} \mathrm{T} (\mathbf {M} ^ {\prime \prime})
\]

是正合的。（§1，第 3 目。）

定义 2. 称右 A-模 E 为平坦的，如果它具有命题 1 的诸等价性质。

左 A-模的平坦性可类似地定义。说一个右 A-模 E 是平坦的，等价于说 E 视为左 A \(^{0}\) -模时是平坦的。

注记

\begin{enumerate}
\def\labelenumi{(\arabic{enumi})}
\item
  由第 2 目的引理 3，为使右 A-模 E 是平坦的，必要且充分的是：对 A 的每个有限生成左理想 a，典范映射 \(E \otimes_{A} a \to E\) （命题 1，其像为 Ea）是单射的。
\item
  设 E 是平坦右 A-模。若 M' 是左 A-模 M 的子模，则典范单射 \(E \otimes_{A} M' \to E \otimes_{A} M\) 使我们可以把 \(E \otimes_{A} M'\) 等同于 \(E \otimes_{A} M\) 的一个子群。在此条件下，设 N 是一个左 A-模， \(u: M \to N\) 是一个同态，K = Ker u，I = Im u。考虑正合列
\end{enumerate}

\[
0 \longrightarrow \mathrm{K} \longrightarrow \mathrm{M} \stackrel {u} {\longrightarrow} \mathrm{N}
\]

即可看出（命题 1）， \(\mathbf{E} \otimes_{\mathbf{A}} (\operatorname{Ker} u)\) 等同于 \(\operatorname{Ker}(1_{\mathbf{E}} \otimes u)\) 。另一方面，用 \(u'\) 表示满同态 \(\mathbf{M} \to \mathbf{I}\) ，它与 \(u\) 有相同的图，用 \(i\) 表示典范单射 \(\mathbf{I} \to \mathbf{N}\) ，则 \(1_{\mathbf{E}} \otimes u'\) 是满射的（第 1 目，引理 1），且 \(1_{\mathbf{E}} \otimes i\) 是单射的（因 \(\mathbf{E}\) 是平坦的）。由于

\[
\mathbf {1} _ {\mathbf {E}} \otimes \boldsymbol {u} = (\mathbf {1} _ {\mathbf {E}} \otimes i) \circ (\mathbf {1} _ {\mathbf {E}} \otimes u ^ {\prime}),
\]

\(E \otimes_{A} (Im u) \text{ 等同于 } Im(1_{E} \otimes u).\)

命题 2. (i) 设 \((\mathbf{E}_t)_{t \in I}\) 是一族右 A-模。为使 \(\mathbf{E} = \bigoplus_{t \in I} \mathbf{E}_t\) 是平坦的，必要且充分的是诸 \(\mathbf{E}_t\) 都是平坦的。

\begin{enumerate}
\def\labelenumi{(\roman{enumi})}
\setcounter{enumi}{1}
\tightlist
\item
  设 \(\mathbf{I}\) 是一个有序集， \((\mathrm{E}_{\alpha}, f_{\beta \alpha})\) 是右 A-模的一个正向系（《代数学》第二章 §6 第 2 目）。若诸 \(\mathrm{E}_{\alpha}\) 中的每一个都是平坦的，则 \(\mathbf{E} = \varinjlim \mathrm{E}_{\alpha}\) 是平坦的。
\end{enumerate}

设 \(M' \to M\) 是一个单射的左 A-模同态。

\begin{enumerate}
\def\labelenumi{(\roman{enumi})}
\tightlist
\item
  为使直和同态
\end{enumerate}

\[
\bigoplus_ {i \in I} \left(E _ {i} \otimes_ {A} M ^ {\prime}\right)\rightarrow \bigoplus_ {i \in I} \left(E _ {i} \otimes_ {A} M\right)
\]

是正合的，必要且充分的是诸同态 \(\mathbf{E}_{\iota} \otimes_{\mathbf{A}} \mathbf{M}' \to \mathbf{E}_{\iota} \otimes_{\mathbf{A}} \mathbf{M}\) 中的每一个都如此（《代数学》第二章 §1 第 6 目命题 7 的推论 1）；由于 \(\bigoplus_{\iota \in \mathbf{I}} (\mathbf{E}_{\iota} \otimes_{\mathbf{A}} \mathbf{M})\) 典范地等同于 \(\mathbf{E} \otimes_{\mathbf{A}} \mathbf{M}\) （《代数学》第二章 §3 第 7 目命题 7），这就证明了 (i)。

\begin{enumerate}
\def\labelenumi{(\roman{enumi})}
\setcounter{enumi}{1}
\tightlist
\item
  由假设，诸序列
\end{enumerate}

\[
0 \rightarrow \mathrm{E} _ {\alpha} \otimes_ {\mathrm{A}} \mathrm{M} ^ {\prime} \rightarrow \mathrm{E} _ {\alpha} \otimes_ {\mathrm{A}} \mathrm{M}
\]

都是正合的，因而序列

\[
0 \rightarrow \mathrm{E} \otimes_ {\mathrm{A}} \mathrm{M} ^ {\prime} \rightarrow \mathrm{E} \otimes_ {\mathrm{A}} \mathrm{M}
\]

也是正合的，因为取正向极限与张量积可交换（《代数学》第二章 §6 第 3 目命题 7），且保持正合性（出处同上，§6 第 2 目命题 3）。

